
\section{Resultados e Discussão}


Este capítulo apresenta e discute os resultados obtidos
nos experimentos de regressão simbólica realizados com
o PySR, utilizando como referência os preços calculados
pela formulação analítica de Black--Scholes.

A análise considera a precisão preditiva, a complexidade
das expressões identificadas e sua correspondência com
as propriedades matemáticas do modelo original.


\subsection{Resultados dos Experimentos Univariados}
\label{sec:resultados_univariados}

Nesta etapa, são analisados os resultados obtidos pela
variação individual dos parâmetros de Black--Scholes,
mantendo-se os demais constantes.

As expressões identificadas pelo PySR são comparadas
com as expansões analíticas apresentadas na fundamentação
teórica, buscando investigar a recuperação de estruturas
matemáticas e a precisão das aproximações.


\subsection{Tempo Restante até o Vencimento ($\tau$)}
\label{subsec:resultados_tau}

% 1. Apresentação da expressão selecionada pelo PySR
% 2. Equações 6, 14 e 16
% 3. Comparação com a expansão de Puiseux
% 4. Tabela dos coeficientes teóricos e encontrados
% 5. Comparação dos erros e complexidades
% 6. Discussão dos resultados


\subsection{Taxa Livre de Risco ($r$)}
\label{subsec:resultados_r}

% 1. Expressão selecionada pelo PySR
% 2. Equações 8, 10, 14 e demais relevantes
% 3. Comparação com a expansão de Taylor
% 4. Tabela dos coeficientes
% 5. MAE, MSE, RMSE e complexidade
% 6. Discussão da recuperação estrutural


\subsection{Volatilidade ($\sigma$)}
\label{subsec:resultados_sigma}

% 1. Expressão selecionada pelo PySR
% 2. Equações 5, 14 e 15
% 3. Comparação com a expansão de Taylor
% 4. Análise dos coeficientes teóricos
% 5. Comparação preditiva
% 6. Análise das sensibilidades Vega e Vomma


\subsection{Preço do Ativo Subjacente ($S$)}
\label{subsec:resultados_S}

% 1. Expressão selecionada pelo PySR
% 2. Equações relevantes
% 3. Comparação com a expansão de Taylor
% 4. Análise dos coeficientes teóricos
% 5. Avaliação preditiva
% 6. Análise das sensibilidades Delta e Gamma


\section{Discussão Comparativa dos Experimentos Univariados}
\label{sec:discussao_univariados}

% Comparação da recuperação estrutural
% Relação entre complexidade simbólica e erro
% Diferenças entre ajuste local e extrapolação
% Limitações das expressões encontradas


\section{Experimentos Multivariados}
\label{sec:resultados_multivariados}

% Inserir os resultados dos experimentos
% com mais de uma variável independente.
% Preservar aqui as análises com trajetórias
% de Monte Carlo, quando forem discutidas.


\section{Síntese dos Resultados e Limitações}
\label{sec:sintese_resultados}

% Principais resultados
% Limitações metodológicas
% Alcance das conclusões
